\documentclass[../../main.tex]{subfiles}
\begin{document}

{\bf [解]}
$h(x) = f(x) - g(x)$ とおくと，題意より $h(x)$ は3次以下の関数で，
\begin{align}
\begin{cases}
h(0)=h(2)=h(3)=0 \\
h'''(x) = 6
\end{cases} \label{eq:1}
\end{align}
を満たす．$h(x)$が２次以下なら$h'''(x)=0$となるから，$h(x)$ は3次関数である．
これと\cref{eq:1}の一つ目の式および因数定理から$h(x)=0$ は $x=0,2,3$ を解に持つ．
$h(x)=0$ は3次方程式ゆえ高々3つの解しか持たないこととあわせて，
\begin{align}
h(x) = Ax(x-2)(x-3)
\end{align}
と表せる．ここで $A\in\mathbb{R}\neq 0$である．
この時 $h'''(x)=6A$ だから\cref{eq:1}の第二式に代入して $A=1$を得るから，
\begin{align}
h(x) = x(x-2)(x-3) \label{eq:2}
\end{align}
である．

そこで以下題意の$F(\theta)$を考えよう．簡単のため
$P=\displaystyle\int_1^3 f(x)dx$, $Q=\displaystyle\int_1^3 g(x)dx$ とおくと
\begin{align}
P-Q = \int_1^3 h(x)dx = \left[\dfrac{1}{4}x^4-\dfrac{5}{3}x^3+3x^2\right]_1^3 = \dfrac{7}{3} \quad \cdots \text{③} \label{eq:3}
\end{align}
のもとで，
\begin{align}
F(\theta) = Pc^2+Qs^2 \quad (c=\cos\theta,\ s=\sin\theta) \label{eq:4}
\end{align}
が $0 \le\theta\le 2\pi$ で最大値をとる $\theta$ をもとめれば良い．
\cref{eq:3}から $Q=P-\dfrac{7}{3}$ だから\cref{eq:4}に代入して
\begin{align}
F(\theta) = P(1-s^2)+\left(P-\dfrac{7}{3}\right)s^2 = -\dfrac{7}{3}s^2+P
\end{align}
である．$0\le s^2\le 1$ よりこれは$s^2=0$の時に最大値$P$をとる．
よって求める$\theta$の値は $\theta=0,\pi,2\pi$ である．

\newpage

\end{document}
